\chapter{热平衡态下微观状态的统计分布律}
\problem{2.1}{有 $N$ 个粒子，当 $0<v\leq v_0$ 时，其速率分布函数为 $f(v)=C$；当 $v>v_0$ 时，$f(v)=0$；其中 $v_0$ 为常量。
\begin{enumerate}
\item 试作出速度分布曲线；
\item 试确定常数 $C$ 和粒子的平均速率。
\end{enumerate}}
\answerarea
\problem{2.2}{已知速率分布函数为 $f(v)$，试确定速率在 $v_1\sim v_2$ 区间内的分子之平均速率。}
\answerarea
\clearpage
\problem{2.3}{试根据麦克斯韦速率分布律，确定速率的倒数的平均值。}
\answerarea
\problem{2.4}{试根据麦克斯韦分布律证明：分子平动动能在 $\varepsilon$ 到 $\varepsilon+\mathrm d\varepsilon$ 区间的概率为
\[
f(\varepsilon)\,\mathrm d\varepsilon
=\frac{2}{\sqrt{\pi}}(k_{\rm B}T)^{-3/2}
e^{-\varepsilon/(k_{\rm B}T)}\sqrt{\varepsilon}\,\mathrm d\varepsilon,
\]
其中 $\varepsilon=\frac12mv^2$，并根据上式确定分子平动动能的最概然值。}
\answerarea
\clearpage
\problem{2.5}{气体分子局限于一维运动，速率服从麦克斯韦分布律，试确定这些气体分子的方均根速率、平均速率和最概然速率。}
\answerarea
\problem{2.6}{气体分子局限于二维运动，速度的每个分量都服从麦克斯韦分布律，试确定这些气体分子的方均根速率、平均速率和最概然速率。}
\answerarea
\clearpage
\problem{2.7}{系统 1 和 2 都是满足麦克斯韦速率分布律的理想气体系统，两系统中分子的最概然速率分别是 $v_{p1}$ 及 $v_{p2}$。试证明：
\[
\frac{F_{\rm M}(v_{p1})}{F_{\rm M}(v_{p2})}=\frac{v_{p1}}{v_{p2}}.
\]}
\answerarea
\problem{2.8}{某种气体分子在温度 $T_1$ 时的方均根速率等于温度为 $T_2$ 时的平均速率，试确定 $T_2/T_1$。}
\answerarea
\clearpage
\problem{2.9}{试计算 $300\,\mathrm K$ 温度下，氧分子的最概然速率、平均速率和方均根速率。}
\answerarea
\problem{2.10}{穿过某原子能反应堆中心处一面积的中子流密度是 $4\times10^{16}/(\mathrm{m^2\cdot s})$，假设这些中子是温度为 $300\,\mathrm K$ 的“热”中子，并服从麦克斯韦速率分布律，试确定中子气的数密度及分压强。}
\answerarea
\clearpage
\problem{2.11}{试根据麦克斯韦速度分布律，确定气体分子速度分量 $v_x$ 的平方的平均值，并由此给出气体分子每一平动自由度所具有的平动能。}
\answerarea
\problem{2.12}{假设地球表面附近的大气为温度为 $0\,^\circ\mathrm C$、压强为 $1\,\mathrm{atm}$ 的理想气体，速率大于地球上的第二宇宙速率的分子到达大气层边界就逃离大气层消失于星际空间中。
\begin{enumerate}
\item 对平均摩尔质量为 $29\times10^{-3}\,\mathrm{kg/mol}$ 的空气，试确定在一秒钟时间内通过单位面积和地球表面空气中损失的分子数（泄漏数率），以及 1 万年时间内地球上总共损失的空气的质量；
\item 如果前述大气为氢气，试确定其泄漏数率；
\item 假定在地球诞生以来的约 50 亿年内，其表面附近的大气的温度和压强没有发生过变化，试通过具体计算说明现在地球附近的大气中没有氢气。
\end{enumerate}}
\answerarea
\clearpage
\problem{2.13}{观测测得月球的质量是地球的约 $1.2\%$、半径是地球的约 $27\%$，其上的重力加速度是地球上的约 $1/6$，其表面的日平均气温大约为 $-50\,^\circ\mathrm C$（白天受阳光直射处的温度可高达 $127\,^\circ\mathrm C$，夜晚其温度可降至 $-183\,^\circ\mathrm C$），试通过计算具体说明月球上仅可能有极稀薄的且结构复杂的大气。}
\answerarea
\problem{2.14}{试确定速率大于某一 $v_0$ 的气体分子每秒时间内与单位面积器壁的碰撞次数。}
\answerarea
\clearpage
\problem{2.15}{通过测量从容器薄壁上小孔在 1 秒钟内流出的气体的质量 $\omega$，即可确定其中很稀薄的气体的压强。如果小孔的面积为 $S$，则有：
\[
p=\frac{\omega\sqrt{2\pi RT}}{S\sqrt{\mu}},
\]
其中 $\mu$ 是气体的摩尔质量，容器外被抽成真空。
\begin{enumerate}
\item 证明该公式；
\item 这一方法为什么只适用于很稀薄的气体和在薄壁上的开孔？
\end{enumerate}}
\answerarea
\problem{2.16}{一容器被薄隔板分成两部分，其中气体的压强、分子数密度分别为 $p_1,n_1$；$p_2,n_2$。两部分气体的温度相同，都等于 $T$；摩尔质量也相同，都为 $\mu$。试证明：如果 $p_1>p_2$，隔板上有一面积为 $A$ 的小孔，则每秒时间内通过小孔由压强为 $p_1$ 的一方净流入压强为 $p_2$ 的一方的气体质量为
\[
M=\sqrt{\frac{\mu}{2\pi RT}}A(p_1-p_2).
\]}
\answerarea
\clearpage
\problem{2.17}{天然铀矿石的成分主要是稳定的 $^{238}\mathrm U$，仅有少量放射性核素 $^{235}\mathrm U$。为生产需要，希望由之生产出放射性核素丰度很高的核燃料。生产中，常采用先制备出 $\mathrm{UF_6}$ 气体，然后通过多级飘泻流的方法实现对 $^{235}\mathrm U$ 的提纯。现有 $\mathrm{UF_6}$ 气体中 $^{235}\mathrm U$ 的丰度为 $0.7\%$，$^{238}\mathrm U$ 的丰度为 $99.3\%$。为得到 $^{235}\mathrm U$ 丰度为 $99.0\%$ 的核燃料，需要进行多少级飘泻流？}
\answerarea
\problem{2.18}{设有一体积为 $V$ 的容器置于真空中，容器内充满温度恒定为 $T$ 的稀薄气体，容器壁上开有一个面积为 $\Delta S$ 的小孔，若容器壁的厚度可以忽略不计。
\begin{enumerate}
\item 试导出单位时间内由小孔射出的气体分子的数目与气体的压强 $p$ 及温度 $T$ 的关系；
\item 试求出容器中的气体的压强减少到其初始压强的一半所需要的时间（用 $V,\Delta S$ 和容器中气体分子的平均速率 $\overline v$ 表示）。
\end{enumerate}}
\answerarea
\clearpage
\problem{2.19}{二维电子气是目前描述电子器件中与电现象相关问题的常用模型。现有一器件，其表面呈边长为 $L$ 的正方形，其中电子的数目记为 $N$，每个电子的质量记为 $m$。
\begin{enumerate}
\item 记这部分电子气中电子的面密度为 $n$，电子的平均动能为 $\overline\varepsilon$，试证明该电子气系统的压强可以近似表述为 $p=n\overline\varepsilon$。
\item 试求出从边缘上线段为 $\Delta L$ 的狭缝漏流出的电子的速度分布律。
\item 记每个电子携带的电量为 $-e$，试确定这些漏流出的电子可以形成的电流的电流强度。
\end{enumerate}}
\answerarea
\problem{2.20}{根据低温下稀薄气体的容器壁上可吸附气体分子的现象，人们发明了可以提高真空度的“低温泵”。考虑一半径为 $0.2\,\mathrm m$ 的球形容器，除对其壁上一面积为 $1.0\,\mathrm{cm^2}$ 的区域冷却到 $77\,\mathrm K$ 外，其余部分及其内部都保持温度 $300\,\mathrm K$。假设初始时刻容器中水蒸气的压强为 $1.33\,\mathrm{Pa}$，每个水分子碰到上述被冷却的区域即被吸附或凝结到上面，试问需要多长时间可以使容器内的压强降到 $1.33\times10^{-4}\,\mathrm{Pa}$？}
\answerarea
\clearpage
\problem{2.21}{设地面附近的大气是等温的，温度为 $t=5.0\,^\circ\mathrm C$，海平面上的气压为 $p_0=750\,\mathrm{mmHg}$，今测得某山顶的气压 $p=590\,\mathrm{mmHg}$，试确定该山的高度。已知空气的平均摩尔质量为 $28.97\times10^{-3}\,\mathrm{kg/mol}$。}
\answerarea
\problem{2.22}{假定在海平面上大气压是 $1.00\times10^5\,\mathrm{Pa}$，那么，在一般喷气式飞机飞行的海拔 $1.00\times10^4\,\mathrm m$ 高空处大气压是多少？珠穆朗玛峰顶海拔 $8882\,\mathrm m$，那里的大气压是多少？设大气层等温，保持为 $0\,^\circ\mathrm C$。}
\answerarea
\clearpage
\problem{2.23}{试确定在等温大气中一个分子的平均重力势能。}
\answerarea
\problem{2.24}{飞机起飞前，舱中压力计指示为 $1.0\,\mathrm{atm}$，温度为 $27\,^\circ\mathrm C$；起飞后，压力计指示为 $0.80\,\mathrm{atm}$，温度仍为 $27\,^\circ\mathrm C$，试计算飞机距地面的高度。（空气的平均摩尔质量是 $28.97\times10^{-3}\,\mathrm{kg/mol}$）}
\answerarea
\clearpage
\problem{2.25}{实验测得地球表面的大气压强为 $1.0\,\mathrm{atm}$，温度为 $0.0\,^\circ\mathrm C$，地球的半径为 $R_0=6370\,\mathrm{km}$，重力加速度 $g=9.8\,\mathrm{m/s^2}$ 保持不变，大气的平均摩尔质量为 $\mu=28.9\times10^{-3}\,\mathrm{kg/mol}$，试在等温大气模型下计算围绕地球的大气的总质量。}
\answerarea
\problem{2.26}{实际测量表明，距离地面不太远的大气层的温度随高度增加而近似线性降低，即有 $T(z)=T_0-\alpha z$，其中 $\alpha$ 为常量，$T_0$ 为地球表面的温度。记地球表面大气的压强为 $p_0$，大气的平均摩尔质量为 $\mu$，大气层中的重力加速度 $g$ 为常量，试给出高度为 $z$ 处的大气压强。}
\answerarea
\clearpage
\problem{2.27}{试确定温度为 $27\,^\circ\mathrm C$ 时，$1\,\mathrm{mol}$ 氧气具有的平动能和转动能各为多少？}
\answerarea
\problem{2.28}{试确定在室温 $300\,\mathrm K$ 下，$1\,\mathrm{mol}$ 氢和 $1\,\mathrm{mol}$ 氦的内能各是多少？1 克氢和 1 克氦呢？}
\answerarea
\clearpage
\problem{2.29}{试确定水蒸气分解为同温度的氢气和氧气时，内能增加了百分之几？}
\answerarea
\problem{2.30}{试确定在 $3\,\mathrm{kg/cm^2}$ 的压强下，5 升氢气的热运动能量。}
\answerarea
\clearpage
\problem{2.31}{试确定常温下质量为 $M_1=3.00\,\mathrm g$ 的水蒸气与 $M_2=3.00\,\mathrm g$ 的氢气的混合气体的定容比热。}
\answerarea
\problem{2.32}{容器中盛有汞蒸气（$\mathrm{Hg}$）、氖（$\mathrm{Ne}$）、氦（$\mathrm{He}$）的混合气体，试比较温度相同情况下这三类分子的平均动能和方均根速率。}
\answerarea
\clearpage
\problem{2.33}{对于欧洲核子中心的 LHC 的高能铅—铅碰撞实验，假设碰撞后很短时间内形成的“火球”内的物质是由质量很小、速度很接近光速的夸克形成的物质，试估算这些夸克物质的定压比热容。}
\answerarea
\problem{2.34}{假设粒子按微观态的分布相对玻尔兹曼分布的偏离为 $\Delta N_i/N_{i,\rm B}\approx10^{-8}$，试计算系统的微观态数目的相对偏离 $(\Omega_{\rm B}+\Delta\Omega)/\Omega_{\rm B}$。}
\answerarea
\clearpage
\problem{2.35}{暴露在分子质量为 $m$、分子数密度为 $n$、温度为 $T$ 的理想气体中的干净的固体表面可以以某一数率吸收气体分子。如果固体仅完全吸收撞击到其表面上的速度法向分量至少为 $v_s$ 的气体分子，而完全不吸收速度法向分量小于 $v_s$ 的气体分子，试给出固体对气体分子的吸收数率。}
\answerarea
\problem{2.36}{对处于容积 $V$ 内的由 $N$ 个单原子分子组成的气体加热到足够高的温度，则这些气体原子中的一部分可以失去一个电子而形成离子，这种现象称为热电离。在温度为 $T$ 的情况下，热电离达到平衡时，气体中未电离的原子、离子和自由电子的数密度分别为 $n_a,n_i,n_e$，记原子的电离能为 $E_i$，一个原子所占的体积为 $V_a$，试证明：
\begin{enumerate}
\item 上述各量之间满足关系：
\[
n_i=n_e,\quad n_i+n_a=\frac NV,\quad
V_a=\frac{n_a}{n_i n_e}e^{-E_i/(k_{\rm B}T)}；
\]
\item 记这些气体分子的电离度为 $f$，则它与前述各量间满足关系：
\[
\frac{f^2}{1-f}\frac NV=\frac1{V_a}e^{-E_i/(k_{\rm B}T)}.
\]
\end{enumerate}}
\answerarea
\clearpage
\problem{2.37}{设一系统有 4000000 个分子，分子只具有三个可能的分立状态，其能量分别为 $0,\varepsilon,2\varepsilon$，角动量量子数分别为 $0,1,2$。初始时，在 0 能级上有 2000000 个分子，在 $\varepsilon$ 能级上有 1600000 个分子，在 $2\varepsilon$ 能级上有 400000 个分子。如果分子由高能量态退激发到低能量态，则会发出光，并且光子的角动量量子数为 1。
\begin{enumerate}
\item 试通过定量计算说明，对系统进行光谱测量时，是否可以测量到发射或/和吸收谱；
\item 若既能测量到发射谱又能测量到吸收谱，在不考虑能量等对跃迁强度直接影响的情况下，发射谱和吸收谱的强度比是多少？
\end{enumerate}}
\answerarea
\problem{2.38}{通过测量星体表面附近的摩尔质量为 $\mu$ 的气体中原子的光谱线的强度因多普勒红移的变化量和红移量可以确定星体表面的温度。如果测得某一频率 $\nu_0$ 的光谱线的强度减弱一半时的频率红移量为 $\Delta\nu$，试证明星体表面的温度大约为
\[
T=\frac{\mu c}{8R\ln2}\left(\frac{\Delta\nu}{\nu_0}\right)^2.
\]}
\answerarea
\clearpage
\problem{2.39}{现在的测量表明，地球绕太阳的运动可以近似为圆周运动，周期是 365.25 天，太阳表面的温度约为 $5800\,\mathrm K$，它由太阳内部氢原子核聚变为氦原子核的核反应释放出的能量来维持，太阳的质量为 $1.99\times10^{30}\,\mathrm{kg}$，太阳的半径为 $6.96\times10^8\,\mathrm m$，地球到太阳的平均距离为 $1.50\times10^{11}\,\mathrm m$，火星到太阳的平均距离为 $2.28\times10^{11}\,\mathrm m$。
\begin{enumerate}
\item 假设太阳常数（太阳照射到地面上每平方米面积的功率）为 $1.36\,\mathrm{kW/m^2}$，试确定火星表面的平均温度和月亮面向太阳的一面的平均温度。
\item 实验还测得基本电荷电量为 $1.60\times10^{-19}\,\mathrm C$，质子的质量为 $1.67\times10^{-27}\,\mathrm{kg}$，每燃烧 4 个氢核可获得约 $25\,\mathrm{MeV}$ 的热能，太阳质量中的约 $70\%$ 是氢，其中约 $10\%$ 可以燃烧，试估计燃烧这些氢能维持太阳辐射多长时间。
\item 现在的研究表明，地球的年龄大约为 50 亿年，试确定地球诞生时，其绕太阳运动的周期与现在的周期之差。
\item 试确定地球诞生时的太阳常数与现在的太阳常数之差。
\end{enumerate}}
\answerarea
\problem{2.40}{试估算一位身高为 $1.80\,\mathrm m$ 的人（正常体表温度大约 $35.5\,^\circ\mathrm C$）在正常的舒适环境（温度大约为 $20\,^\circ\mathrm C$）下的辐射功率。}
\answerarea
\clearpage
\problem{2.41}{现在很多机场、车站、码头都设有非接触型的红外体温计，以检测是否有高烧的病人。
\begin{enumerate}
\item 如果认定体温在 $39.5\,^\circ\mathrm C$（忽略衣服表面温度与人体温度的差异）及其以上的人为高烧病人，试确定光敏体温计前滤光片最敏感的红外光波长的上限；
\item 如果光敏体温计也能同样检测到正常体温（记之为 $35.5\,^\circ\mathrm C$）的人的频谱，试确定该体温计检测到的体温为 $39.5\,^\circ\mathrm C$ 的高烧病人的频谱亮度比其检测正常体温的人的频谱亮度高出的百分率。
\end{enumerate}}
\answerarea
\problem{2.42}{电子管的真空度为 $1.333\times10^{-3}\,\mathrm{Pa}$，设空气分子的有效直径为 $3.0\times10^{-10}\,\mathrm m$，试求 $27\,^\circ\mathrm C$ 时单位体积内的分子数 $n$、平均自由程 $\overline\lambda$ 和平均碰撞频率 $\overline Z$。}
\answerarea
\clearpage
\problem{2.43}{在 $1.0\,\mathrm{atm}$、$15\,^\circ\mathrm C$ 情况下，氢气分子的平均自由程为 $1.18\times10^{-7}\,\mathrm m$，试确定氢分子的有效直径。}
\answerarea
\problem{2.44}{设某种气体分子的平均自由程为 $\overline\lambda$，试证明：一个分子在连续两次碰撞之间所走路程至少为 $x$ 的概率为 $e^{-x/\overline\lambda}$。}
\answerarea
\clearpage
\problem{2.45}{假设某压强下温度为 $300\,\mathrm K$ 的氧气的分子都以平均速率运动，并且其平均自由程为 $2.0\,\mathrm{cm}$。某一时刻这些氧气中有 $N$ 个分子都刚与其它分子碰撞，试确定经过多长时间后其中尚有一半未与其它分子相碰。}
\answerarea
\problem{2.46}{某种有效直径为 $2.6\times10^{-10}\,\mathrm m$ 的气体分子在 $25\,^\circ\mathrm C$ 情况下的平均自由程为 $2.63\times10^{-7}\,\mathrm m$，试确定：
\begin{enumerate}
\item 这种情况下气体的压强；
\item 一个分子在 $1.0\,\mathrm m$ 的路程上与其它分子碰撞的次数。
\end{enumerate}}
\answerarea
\clearpage
\problem{2.47}{假设温度为 $T$ 的稀薄混合气体由有效直径分别为 $d_1,d_2$、质量分别为 $m_1,m_2$ 的两类分子组成，并且这两类分子的数密度分别为 $n_1,n_2$，试证明这两类分子的平均自由程分别为
\[
\begin{aligned}
\overline\lambda_1&=\frac{1}{\sqrt2\pi d_1^2 n_1+\pi(d_1+d_2)^2n_2\sqrt{\dfrac{m_1}{16\mu}}},\\[5pt]
\overline\lambda_2&=\frac{1}{\sqrt2\pi d_2^2 n_2+\pi(d_1+d_2)^2n_1\sqrt{\dfrac{m_2}{16\mu}}},
\end{aligned}
\]
其中 $\mu=\dfrac{m_1m_2}{m_1+m_2}$ 为两类分子的折合质量。}
\answerarea
\problem{2.48}{一显像管的灯丝到荧光屏的距离为 $20\,\mathrm{cm}$，已知空气分子的直径为 $3.0\times10^{-10}\,\mathrm m$，在气体温度为 $27\,^\circ\mathrm C$ 的情况下，为使灯丝发射的电子有 $90\%$ 直接达到荧光屏而在途中不与空气分子相碰，显像管的真空度至少为多高？}
\answerarea
\clearpage
\problem{2.49}{由电子枪发出一束电子射入压强为 $p$ 的气体中，在电子枪前相距 $x$ 处放置一收集电极，用来测定能自由通过这段距离的电子数。已知电子枪发射的电子流强度为 $100\,\mu\mathrm A$，当气压 $p=100\,\mathrm{Pa}$、$x=10\,\mathrm{cm}$ 时，到达收集极的电子流强度为 $37\,\mu\mathrm A$。试求：
\begin{enumerate}
\item 电子的平均自由程；
\item 当气压降到 $50\,\mathrm{Pa}$ 时，到达收集极的电子流强度。
\end{enumerate}}
\answerarea
\problem{2.50}{在带电粒子加速器装置中，产生被加速粒子的离子源与加速腔之间都不可避免地有小的“空隙”。通常情况下，被加速粒子的体积与气体分子体积相比可以忽略不计，被加速粒子运动的速度比气体分子运动的速率大得多。若一加速器中上述空隙的长度为 $0.10\,\mathrm m$，为使进入加速腔的粒子流强度不低于离子源处的 $99.999999\%$，则空隙中的真空度至少为多高？}
\answerarea
\clearpage
